\documentclass[hyperref={pdfpagelabels=false}]{beamer}

% Configurações Básicas --------------------------------------------------------
\let\Tiny=\tiny
\usepackage[utf8]{inputenc}
\usepackage[T1]{fontenc}
\usepackage[brazil]{babel}
\usepackage{graphicx}

% Informações Pessoais ---------------------------------------------------------
\author[ROSA]{Wanderson Henrique Camargo Rosa}
\title[JS]{Javascript}
\subtitle{Breve Descrição e Histórico da Linguagem}
\institute[UNISINOS]{Linguagens de Programação 2011/1\\Centro de
Ciências Exatas e Tecnológicas\\Universidade do Vale do Rio dos Sinos ---
UNISINOS}

% Documento --------------------------------------------------------------------
\begin{document}

% Apresentação -----------------------------------------------------------------
\begin{frame}
    \maketitle{}
\end{frame}

\begin{frame}{Javascript}{Sobre a Linguagem de Programação}
    Javascript\cite{mozilla,wikipedia} é uma linguagem de script orientada a
    objetos baseada em protótipos, dinâmica, com tipagem fraca e funções de
    primeira classe.
\end{frame}

\begin{frame}{Javascript}{Ambientes de Utilização}
    Basicamente, o Javascript é utilizado principalmente em navegadores
    \emph{web}\cite{mozilla} no cliente, porém podem ser executados em
    servidores\cite{mozilla} que recebem e enviam informações Javascript.

    Também pode ser encontrado em documentos no formato PDF, navegadores
    específicos e \emph{widgets} para áreas de trabalho\cite{wikipedia}.
\end{frame}

\begin{frame}{Javascript}{Histórico}
    A linguagem de programação Javascript foi criada por Brendan
    Eich\cite{mozilla,wikipedia} na Netscape, sobre um motor chamado
    SpiderMonkey\cite{mozilla}, implementado em C. Primeiramente nomeado como
    Mocha\cite{wikipedia}, após Livescript\cite{history,diff,wikipedia,plotter}
    e por fim seu nome atual, o Javascript foi idealizado para tornar as páginas
    mais dinâmicas, evitando uma nova requisição para o serviço
    \emph{web}\cite{history}.
\end{frame}

\begin{frame}{Javascript}{Criador da Linguagem}
    \begin{figure}
        \includegraphics[scale=0.35]{brendan-eich}
        \caption{Brendan Eich. Twitter @brendaneich}
    \end{figure}
\end{frame}

\begin{frame}{Javascript}{Padronização Internacional}
    Javascript recebeu em 1996 uma padronização internacional\cite{history} com
    a \emph{European Computer Manufacturers Association}, sob o nome ECMA-262 ou
    ECMAScript. Existem outras linguagens baseadas no Javascript, como a JScript
    da Microsoft\cite{diff}. A empresa afirma que o JScript não é executável
    sobre ambientes ECMAScript\cite{wikipedia}.
\end{frame}

\begin{frame}{Javascript}{Programação Orientada a Protótipos}
    Programação baseada em protótipos\cite{prototype} é um estilo de linguagem
    de programação orientada a objetos em que classes não estão presentes. A
    reutilização de código é dada com a clonagem de objetos existentes. Este
    estilo de programação também pode ser chamado de Programação Orientada a
    Protótipos\cite{prototype}.

    Cada objeto possui uma referência para o objeto clonado e não nula,
    fornecendo o encadeamento de protótipos\cite{ecma}.
\end{frame}

\begin{frame}{Javascript}{Influências da Linguagem}
    Segundo a Wikipédia\cite{wikipedia}, Javascript foi baseado nas linguagens
    Self e Scheme. Porém, o grafo de Histórico das Linguagens de Programação
    desenvolvido por Lévénez\cite{plotter} afirma que a linguagem é influenciada
    diretamente pela Cmm (C Minus Minus\cite{cmm}).
\end{frame}

\begin{frame}{Javascript}{Influências da Linguagem}
    \begin{description}
    \item[2002] ECMAScript Draft 04
    \item[2002] Javascript 2.0 Draft 04
    \item[1999] Javascript 2.0 Draft 01
    \item[1999] Javascript 1.5
    \item[1995] Javascript
        \begin{description}
        \item[1999] ECMAScript
        \item[1996] JScript
        \end{description}
    \item[1995] Livescript
    \item[1992] Cmm
        \begin{description}
        \item[1990] ISO C C90
        \item[1989] ARM C++
        \end{description}
    \end{description}
\end{frame}

\begin{frame}{Javascript}{Linguagens Influenciadas}
    Ainda segundo a Wikipédia\cite{wikipedia}, são linguagens influenciadas pelo
    Javascript:
    \begin{itemize}
        \item JScript
        \item JScript .NET
        \item Objective-J
        \item TIScript
    \end{itemize}
\end{frame}

\begin{frame}
    \bibliographystyle{plain}
    \bibliography{document}
\end{frame}

\end{document}